\chapter{Complete Walkthrough of the Source Code}
\label{ch:source}

This chapter walks the three source files from top to bottom. Every code block below is
\textbf{quoted directly out of the project files with \LaTeX{}'s \code{lstinputlisting}}, so the line
numbers printed in the margin are the file's real line numbers --- open \code{hand_logic.py}, jump to
line 127, and you will be looking at exactly what is printed here. The complete files are reproduced
verbatim in Appendix~\ref{app:source}.

The three files are presented in dependency order: the pure logic first, then the application that
drives it, then the tests that pin both down.

\section{\code{hand\_logic.py} --- the pure layer}
\label{sec:src-logic}

\subsection{Module docstring, imports and constants}

\begin{codecard}
\lstinputlisting[firstline=1, lastline=30, firstnumber=1]{../hand_logic.py}
\end{codecard}

\begin{itemize}
  \item \textbf{Lines 1--9, the docstring.} It states two things a reader cannot infer from the code:
        the \emph{dependency policy} (``imports nothing but the standard library'') and the
        \emph{coordinate convention} (origin top-left, $y$ grows downward). The second is the single
        most consequential fact in the module --- every comparison below depends on it --- so it is
        declared before any code.
  \item \textbf{Line 11, \code{from \_\_future\_\_ import annotations}.} Makes type annotations lazy:
        they are stored as strings rather than evaluated. That matters because it lets a hint like
        \code{list[dict]} be written on Python 3.8 as well as 3.9+, without a runtime cost. It is the
        standard first line of any modern Python module.
  \item \textbf{Lines 13--14, the only two imports in the module.} \code{math} for \code{math.hypot},
        and \code{deque} for the debouncer's sliding window. No \code{cv2}, no \code{mediapipe} ---
        the dependency arrow from Chapter~\ref{ch:architecture} enforced by construction.
  \item \textbf{Lines 16--21, the landmark index constants.} These are \emph{named} indices. Writing
        \code{PINKY\_MCP} rather than \code{17} at the point of use is what allows
        Section~\ref{sec:palm} to be checked against Section~\ref{sec:src-logic} line by line: the
        code and the mathematics use the same names.
  \item \textbf{Line 23, \code{FINGER\_PAIRS}.} A list of (tip, PIP) tuples in the order the state
        list will use: index, middle, ring, pinky. Note that it deliberately \emph{excludes} the
        thumb, which follows a completely different rule.
  \item \textbf{Lines 25--30, the remaining constants.} \code{EXPECTED\_LANDMARKS = 21} is used as a
        validity check. \code{THUMB\_DISTANCE} and \code{THUMB\_X} are the string names of the two
        thumb rules --- string constants rather than booleans, so they read well in a CLI and in an
        error message. \code{NO\_GESTURE} is the placeholder emitted when a pose matches no pattern,
        and --- importantly --- the initial value of the debouncer, which is what defines its
        warm-up behaviour.
\end{itemize}

\subsection{The distance helper}

\begin{codecard}
\lstinputlisting[firstline=33, lastline=34, firstnumber=33]{../hand_logic.py}
\end{codecard}

\begin{itemize}
  \item \code{math.hypot(dx, dy)} computes $\sqrt{dx^2 + dy^2}$ without forming the squares
        explicitly, and without the overflow that naive squaring can cause for very large values.
        Here the coordinates are pixels, so overflow is not a concern --- but the substitution also
        avoids a temporary and reads exactly like the formula in Chapter~\ref{ch:math}.
  \item The leading underscore is a Python convention: \code{\_distance} is \emph{internal}. It is
        not imported by any test, signalling that callers outside the module should not depend on it.
\end{itemize}

\subsection{Normalized to pixel conversion, and the bounding box}

\begin{codecard}
\lstinputlisting[firstline=37, lastline=46, firstnumber=37]{../hand_logic.py}
\end{codecard}

\begin{itemize}
  \item \textbf{Line 39} is the comprehension from Chapter~\ref{ch:foundations}: for each normalized
        landmark, multiply $x$ by the frame width and $y$ by the frame height, and truncate to an
        integer with \code{int()}. Note the argument order \code{(x, y)} in the output tuple ---
        pixels are stored as $(x, y)$ pairs even though the NumPy array is indexed \code{[y, x]}.
        This module therefore speaks ``$(x, y)$'' consistently, and only the array indexing code in
        the application has to remember the difference.
  \item \textbf{Line 44--45} build two parallel lists of $x$ and $y$ coordinates. This is the
        standard idiom for reducing a point list to axis extents, and it is slightly clearer than a
        running min/max loop.
  \item \textbf{Line 46} returns the box in the conventional order
        $(x_{\min}, y_{\min}, x_{\max}, y_{\max})$. The order is fixed and documented, because a
        rectangle whose coordinates are swapped is a silent, plausible-looking bug.
\end{itemize}

\subsection{The palm normal}

\begin{codecard}
\lstinputlisting[firstline=49, lastline=61, firstnumber=49]{../hand_logic.py}
\end{codecard}

\begin{itemize}
  \item \textbf{Line 50--56, the docstring.} It writes the formula out in the comment as well as the
        code \emph{and} states the sign convention. For a function whose output is a bare float whose
        meaning is entirely in its sign, that is essential: without the docstring, the returned
        \code{-8775.0} is meaningless.
  \item \textbf{Lines 57--59} destructure the three landmarks used --- wrist, index MCP, pinky MCP
        --- into named local variables. This costs one line each and removes every possibility of
        confusing \code{x5} with \code{x17} while reading the formula.
  \item \textbf{Line 60} is the cross product's $z$-component exactly as derived in
        Section~\ref{sec:palm}. The parenthesisation is explicit rather than relying on operator
        precedence, because a subtle precedence error here would produce a plausible number with the
        wrong sign.
  \item \code{float(...)} on the return is defensive: heights and widths may arrive as NumPy
        integers, and \code{np.int64 * np.int64} can overflow at 32 bits. Converting to Python's
        arbitrary-precision float sidesteps that entirely.
\end{itemize}

\begin{codecard}
\lstinputlisting[firstline=64, lastline=66, firstnumber=64]{../hand_logic.py}
\end{codecard}

A three-line function that exists to give a sign a name. It is separate from \code{palm\_normal\_z}
so that the \emph{convention} lives in exactly one place: if the camera were ever un-mirrored, this
single line is what would change.

\subsection{Reading the handedness label safely}

\begin{codecard}
\lstinputlisting[firstline=69, lastline=78, firstnumber=69]{../hand_logic.py}
\end{codecard}

\begin{itemize}
  \item \textbf{Lines 70--73}: the attribute chain \code{handedness.classification[0].label} can fail
        in three distinct ways --- the object may lack the attribute, the list may be empty, or the
        value may not be subscriptable at all. Catching all three exception types in one clause is
        correct here because the recovery is identical for each: the label is unknown.
  \item \textbf{Lines 74--75}: an unexpected string is normalised to \code{"Unknown"} rather than
        passed through. Everything downstream then has exactly three possible labels, which is what
        makes the \code{in ("Left", "Right")} checks elsewhere total.
  \item \textbf{Lines 76--77}: \code{--swap-handedness} is applied \emph{here}, at parse time, so
        that the swapped label is what the rest of the program sees --- including the state keys used
        for smoothing and debouncing. Swapping later (at draw time) would have left the state keyed
        by the original, inconsistent label.
\end{itemize}

\subsection{Parsing a frame: the single-pass cache}

\begin{codecard}
\lstinputlisting[firstline=81, lastline=124, firstnumber=81]{../hand_logic.py}
\end{codecard}

This is the most important function in the module, so it is worth reading line by line.

\begin{itemize}
  \item \textbf{Lines 82--92, the docstring.} It states the duck-typing contract explicitly
        (``only ever touches \code{multi\_hand\_landmarks} and \code{multi\_handedness}''), which is
        the reason the tests can pass a fake object and get real coverage. It also documents the
        exact dictionary schema --- six keys --- and the two invariants: never raises, and skips
        incomplete hands.
  \item \textbf{Lines 93--95}, \code{getattr(results, "multi\_hand\_landmarks", None)} and the
        truthiness check. A result with no hands may or may not have the attribute; either way,
        \code{None} or an empty list returns \code{[]}. This one \code{if} is what makes the
        ``zero hands'' and ``\code{None} input'' tests pass without a special case.
  \item \textbf{Lines 97--99}: \code{or []} converts a present-but-\code{None} attribute into an
        empty list, so the length check on the next line is always safe.
  \item \textbf{Lines 101--104, the loop}. \code{enumerate} supplies both the index (used to look up
        the matching handedness entry) and the hand itself. The index correspondence matters:
        MediaPipe emits the two lists in the same order, which is the only reason
        \code{handedness[index]} is correct.
  \item \textbf{Lines 105--107, the completeness gate.} A hand without a full set of 21 landmarks is
        skipped with \code{continue} rather than \code{break}: the \emph{other} hand in the frame may
        be perfectly usable, and discarding it would be a worse outcome than ignoring the malformed
        one.
  \item \textbf{Lines 109--110}: \code{norm} keeps the $(x, y, z)$ triples, and \code{px} converts
        with the same helper described above. Both are stored because they serve different consumers
        --- \code{norm} for debugging and future 3D work, \code{px} for every rule in the module.
  \item \textbf{Lines 111--124, the dictionary literal.} Every derived value is computed \emph{here
        and once}: the palm normal, the orientation label, and the bounding box. This is the
        single-pass cache of Chapter~\ref{ch:architecture} --- the reason the render loop performs no
        redundant coordinate arithmetic.
\end{itemize}

\subsection{The finger rules}

\begin{codecard}
\lstinputlisting[firstline=127, lastline=162, firstnumber=127]{../hand_logic.py}
\end{codecard}

\begin{itemize}
  \item \textbf{Lines 128--139, the docstring.} It documents all three rules --- the vertical test for
        the four long fingers, the distance rule for the thumb, and the alternative horizontal rule
        --- plus the exact safe-result contract. It also explains why \code{palm\_z} is accepted but
        unused by the default rule (API symmetry with the orientation readout), which prevents a
        future reader from assuming the parameter is a bug.
  \item \textbf{Lines 141--142, the input guard.} This single line implements the specification's
        requirement: a \code{None}, empty, or short landmark list returns
        \code{[False] * 5} and never raises. Note it is a \emph{list comprehension-free} early
        return, evaluated before any indexing can occur.
  \item \textbf{Line 146, the four-finger test.} A single comprehension over \code{FINGER\_PAIRS},
        comparing the $y$ of each tip against the $y$ of its PIP joint --- the rule derived in
        Section~\ref{sec:screencs}. The result is a list of exactly four booleans, in the order the
        pairs are declared.
  \item \textbf{Lines 148--150, the horizontal thumb rule.} Two conditions guard it: an explicit
        \code{thumb\_mode == THUMB\_X} \emph{and} a usable label. If either fails, control falls to
        the \code{else} branch --- so an \code{"Unknown"} label degrades to the distance rule instead
        of silently reporting a folded thumb. That fallback is asserted by a test.
  \item \textbf{Lines 151--154, the distance rule.} The default path, and the return of
        Section~\ref{sec:thumb}'s derivation. Written as a direct transcription of the inequality,
        which is why it is readable without a comment.
  \item \textbf{Lines 155--158, the exception guard.} The commented contract is ``never raises'',
        and a comprehension over user-supplied data cannot guarantee that. Catching
        \code{IndexError}, \code{TypeError} and \code{ValueError} covers short lists, \code{None}
        elements and non-numeric entries, and all three recover to the same safe value. A test feeds
        a list of 21 \code{None} values specifically to exercise this path.
  \item \textbf{Line 160}: the thumb result is \emph{prepended} to the four finger results, giving
        the canonical \code{[thumb, index, middle, ring, pinky]} order that every consumer ---
        classification, HUD text, and the tests --- relies on.
\end{itemize}

\subsection{Gesture classification}

\begin{codecard}
\lstinputlisting[firstline=165, lastline=182, firstnumber=165]{../hand_logic.py}
\end{codecard}

\begin{itemize}
  \item \textbf{Lines 166--167}: a short or empty state list returns the placeholder rather than
        raising, keeping the module's no-exceptions contract intact.
  \item \textbf{Line 169} unpacks exactly five values and coerces each with \code{bool()}. The
        coercion matters: it means a caller passing \code{1}/\code{0}, or NumPy booleans, gets the
        same behaviour as one passing \code{True}/\code{False}.
  \item \textbf{Lines 171--180, the pattern chain.} Each test is \emph{exhaustive} --- it states
        which fingers must be down as well as which must be up. Without the negated terms, ``Peace''
        would also match an open palm, and the first match would win. The ordering from most to least
        obvious is a readability choice, not a correctness requirement, because the patterns are
        mutually exclusive.
  \item \textbf{Line 182}: anything unmatched returns \code{NO\_GESTURE}. The module never returns
        \code{None} for a gesture, so the HUD can always print something.
\end{itemize}

\subsection{The debouncer}

\begin{codecard}
\lstinputlisting[firstline=185, lastline=207, firstnumber=185]{../hand_logic.py}
\end{codecard}

\begin{itemize}
  \item \textbf{Lines 185--191, the class docstring.} It states the contract in three clauses: the
        window length, the agreement threshold, and --- crucially --- the warm-up behaviour
        (``before any majority has ever formed it returns \code{---}''). That last sentence is the
        answer to a question the code alone leaves ambiguous.
  \item \textbf{Lines 193--198, the constructor.} Three pieces of state per instance:
        \code{window} and \code{min\_agree} as configuration, \code{\_history} as a
        \code{deque(maxlen=window)} that evicts automatically, and \code{\_stable} seeded with the
        placeholder. Seeding with the placeholder \emph{is} the warm-up policy.
  \item \textbf{Lines 200--201, the \code{stable} property.} A read-only view of the last accepted
        label. Exposing it via \code{@property} rather than a plain attribute means callers cannot
        accidentally assign to it --- a small invariant guard.
  \item \textbf{Lines 203--207, \code{update}.} Three statements: append the new observation (the
        \code{deque} silently evicts the oldest), test whether the \emph{new} label alone has reached
        the threshold within the window, and return the current stable label either way. Counting
        only the newest label is correct, and cheaper than building a histogram of all five entries:
        if any label reached the threshold, the newest one is the only candidate that could have just
        crossed it.
\end{itemize}

\subsection{The exponential moving average}

\begin{codecard}
\lstinputlisting[firstline=210, lastline=220, firstnumber=210]{../hand_logic.py}
\end{codecard}

\begin{itemize}
  \item \textbf{Lines 211--214, the docstring}, which writes the recurrence in words and states the
        first-sample rule.
  \item \textbf{Lines 215--216}: the \code{None} special case. Without it, the arithmetic would
        raise, or --- worse, if a caller substituted zero --- the dot would animate in from the
        origin on the first frame.
  \item \textbf{Lines 217--220}: the weighted sum, applied to $x$ and $y$ independently but in one
        expression, returning a plain tuple. Floats are returned, not integers: rounding here would
        accumulate error and defeat the purpose of smoothing. The caller rounds once, at draw time.
\end{itemize}

\section{\code{hand\_tracker.py} --- the application}
\label{sec:src-tracker}

\subsection{Module setup, the import guard, and constants}

\begin{codecard}
\lstinputlisting[firstline=36, lastline=71, firstnumber=36]{../hand_tracker.py}
\end{codecard}

\begin{itemize}
  \item \textbf{Line 36}: the module logger. Every diagnostic in the application goes through it
        rather than \code{print}, so verbosity is controllable from the command line and library
        users are not spammed by a chatty dependency.
  \item \textbf{The \code{warnings.catch\_warnings()} block} around the MediaPipe import suppresses
        the deprecation noise of the legacy API --- scoped to the import, so real warnings elsewhere
        still surface. This is discussed in detail in Section~\ref{sec:deps}.
  \item \textbf{The \code{RuntimeError} guard} is the third dependency trap from
        Chapter~\ref{ch:engine}: it converts an obscure \code{AttributeError} from deep inside a
        library into one actionable sentence naming the fix.
  \item \textbf{The \code{tikz}-free constants block (lines 54--71)} collects every magic number the
        application uses: the window title, the captures directory, the default frame size, the six
        HUD styling values, the two index-tip marker radii and its colour, and the FPS averaging
        window. Collecting them means the look of the overlay can be tuned without hunting through
        the drawing code --- and it makes each one citable in this handbook.
  \item \code{INDEX\_TIP\_COLOR = (0, 165, 255)} deserves a note: this is \textbf{BGR}, not RGB.
        Read as a colour, it is ``no blue, 165 green, 255 red'' --- the orange dot visible in the
        running application. A constant written as \code{(255, 165, 0)} would draw a \emph{blue}
        marker, and it would look deliberate.
\end{itemize}

\subsection{\code{HandTracker}: construction and the detection call}

\begin{codecard}
\lstinputlisting[firstline=77, lastline=122, firstnumber=77]{../hand_tracker.py}
\end{codecard}

\begin{itemize}
  \item \textbf{Lines 77--96, the constructor.} The five parameters mirror the CLI flags, each with
        the default the specification requires. Four of them are configuration; the fifth,
        \code{mode}, is passed straight to \code{static\_image\_mode}. Note the ordering of
        assignments: the three caches are created before the MediaPipe model, so an object is never
        in a half-built state even if model construction raises.
  \item \textbf{Line 90, the model.} This is the only place in the program where MediaPipe is
        instantiated. Isolating construction to one line makes the three thresholds from
        Chapter~\ref{ch:engine} easy to find and tune.
  \item \textbf{Line 91--93, the three caches.} \code{hands\_data} is the per-frame telemetry;
        \code{\_smooth\_tip} holds one EMA state per hand key; \code{\_debounce} holds one
        \code{GestureDebouncer} per hand key. All three are \emph{instance} attributes --- two
        trackers would not share state.
  \item \textbf{Lines 98--114, the detection itself.} The conversion to RGB, the read-only flag, the
        model call, and the restore of the write flag. Setting \code{rgb.flags.writeable = False} is
        a micro-optimisation that tells NumPy the buffer will not be written, letting the model read
        it without defensive copying.
  \item \textbf{Lines 116--119, the cache rebuild.} Note the order: the image is measured
        (\code{frame.shape[:2]} gives \emph{height, width}), and \code{parse\_hand\_data} is called
        with \code{(width, height)} --- the opposite order, because the function's parameters are
        named \code{width, height}. This is the single most error-prone line in the file, and it is
        correct precisely because both names are explicit at the call site.
  \item \textbf{Lines 121--122, the optional mesh.} Drawing is guarded by both the \code{draw} flag
        and the presence of landmarks. \code{mp\_drawing.draw\_landmarks} is handed two styles: a
        green joint style and a near-white connection style, which is what produces the familiar
        MediaPipe skeleton.
\end{itemize}

\subsection{The accessor methods}

\begin{codecard}
\lstinputlisting[firstline=124, lastline=147, firstnumber=124]{../hand_tracker.py}
\end{codecard}

\begin{itemize}
  \item \textbf{\code{hand\_count}} is a one-liner over the cache. It is a method rather than a
        direct attribute read so that the length and the cache can never disagree.
  \item \textbf{\code{hand\_label}} returns \code{"Unknown"} out of range --- never raises, and never
        returns \code{None}, so its result can be placed directly into HUD text.
  \item \textbf{\code{find\_positions}} is the specification's method, preserved exactly. Two things
        about it are deliberate: it returns a \emph{copy} (\code{list(...)}) so a caller cannot
        mutate the cache, and it returns \code{[]} --- not \code{None} --- for a bad index, so the
        finger rules receive a value they already have a guard for. The docstring records the
        vestigial \code{frame} parameter and the size caveat from Section~\ref{sec:cache}.
  \item \textbf{\code{fingers\_up}} is a one-line delegation to \code{fingers\_state}, carrying the
        instance's \code{thumb\_mode} so the CLI flag reaches the rule without the render loop having
        to pass it around.
\end{itemize}

\subsection{The smoothed tip, orientation, and debounced gesture}

\begin{codecard}
\lstinputlisting[firstline=149, lastline=178, firstnumber=149]{../hand_tracker.py}
\end{codecard}

\begin{itemize}
  \item \textbf{Lines 150--152, \code{index\_tip}'s guards.} Out-of-range index and a too-short
        landmark list both return \code{None}. \code{None} is the right sentinel here --- unlike
        \code{find\_positions}, the caller genuinely has nothing to draw --- and the render loop
        checks for it before calling \code{cv2.circle}.
  \item \textbf{Line 157, the state key.} The EMA state is looked up by hand identity, not by list
        position, which is the fix for the cross-contamination trap in Section~\ref{sec:ema}.
  \item \textbf{Lines 158--160}: the smoothing call, the state update, and the cast to integers. Note
        that the \emph{float} value is what gets stored, and the rounding happens only on the way out
        --- so the filter's memory is never quantised.
  \item \textbf{\code{orientation}} returns the cached label, with an em-dash for out-of-range. It
        performs no arithmetic at all: the sign test was done once, in \code{parse\_hand\_data}.
  \item \textbf{Lines 169--178, \code{gesture}.} The debouncer is created lazily on first use for a
        given key and then reused. Lazy creation is what allows the tracker to adapt to hands
        appearing and disappearing mid-session without a separate bookkeeping pass --- and because
        the key is the label, a hand that briefly leaves the frame and returns resumes its own
        gesture history rather than inheriting another hand's.
\end{itemize}

\begin{codecard}
\lstinputlisting[firstline=179, lastline=189, firstnumber=179]{../hand_tracker.py}
\end{codecard}

\begin{itemize}
  \item \textbf{\code{\_state\_key}} is a static method because it needs no instance state --- it is
        a pure function of its arguments, and marking it \code{@staticmethod} says so.
  \item The fallback key, \code{f"index:\{hand\_index\}"}, is deliberately namespaced with a prefix.
        If the label list ever contained the literal string \code{"0"}, the two key spaces could
        collide; the prefix makes that impossible.
  \item \textbf{\code{close}} releases the MediaPipe graph. It exists so the teardown in
        Section~\ref{sec:teardown} can call one method rather than reaching into
        \code{self.hands}, keeping the model's lifecycle inside the class that owns it.
\end{itemize}

\subsection{The HUD: text assembly and drawing}

\begin{codecard}
\lstinputlisting[firstline=192, lastline=230, firstnumber=192]{../hand_tracker.py}
\end{codecard}

\begin{itemize}
  \item \textbf{Lines 192--194, \code{hud\_lines}.} A pure function that assembles the text: FPS,
        total, then one line per hand. Separating assembly from drawing means the text can be tested
        without a frame --- and it is, by \code{test\_hud\_lines\_formatting}.
  \item \textbf{Lines 198--199}: an empty line list returns immediately, so an empty card is never
        drawn.
  \item \textbf{Lines 201--205, measurement.} \code{cv2.getTextSize} returns the pixel extent of a
        string in the chosen font. Measuring \emph{before} drawing is what allows the card to size
        itself to its content, which is the fix for the two-hand overflow trap in
        Section~\ref{sec:hud}.
  \item \textbf{Lines 207--214, the clamped rectangle.} Four values are computed: the top-left
        corner, and the bottom-right corner clamped to the image with \code{min}. The early return on
        \code{x1 <= x0 or y1 <= y0} handles a frame so small that the card would be degenerate ---
        20-pixel-square frames do occur in tests and in thumbnail previews.
  \item \textbf{Lines 216--217, the blend.} The ROI is a NumPy \emph{view}; \code{addWeighted} reads
        it, and the assignment writes the darkened result back into the same memory. The expression
        \code{1.0 - HUD\_ALPHA} keeps the two weights summing to 1, so brightness is scaled rather
        than shifted --- which is why the card looks like tinted glass rather than a black box.
  \item \textbf{Lines 219--230, the text loop.} The baseline advances by \code{line\_height} each
        iteration, starting one pad plus one text height below the card's top edge. White
        (\code{255, 255, 255}) is chosen because the blended background is guaranteed dark; the
        \code{cv2.LINE\_AA} flag enables anti-aliasing so small text stays legible.
\end{itemize}

\subsection{The per-frame loop body}

\begin{codecard}
\lstinputlisting[firstline=233, lastline=254, firstnumber=233]{../hand_tracker.py}
\end{codecard}

This function is the shared heart of both the live loop and \code{--image} mode --- which is what
guarantees the still-image path exercises the real drawing code rather than a copy of it.

\begin{itemize}
  \item \textbf{Line 238, the loop bound.} \code{range(tracker.hand\_count())} makes both hands in
        the frame fall out naturally, and makes the two-hand requirement of the specification a
        non-event.
  \item \textbf{Lines 239--242, the read-and-read-again pattern.} \code{hand\_label} and
        \code{find\_positions} both come from the cache populated by the preceding
        \code{find\_hands}. No landmark is recomputed below this line.
  \item \textbf{Lines 243--244}: \code{sum(states)} counts the \code{True} values directly. Since the
        booleans are Python \code{bool}, which is a subclass of \code{int}, this is both correct and
        idiomatic.
  \item \textbf{Line 245, the HUD row.} An f-string with four fields: label, count, debounced
        gesture, and the palm orientation in square brackets. This single line is the ``diagnostic
        visual telemetry'' the specification asks for, and every field is traceable to a rule
        derived in Chapter~\ref{ch:math}.
  \item \textbf{Lines 247--251, the marker.} Two filled circles: a white halo of radius 14, then the
        orange dot of radius 10 on top. Drawing the halo first and the dot second is what produces a
        ringed marker that stays visible against both light and dark backgrounds --- the same design
        principle as the HUD card.
\end{itemize}

\subsection{Screenshots, and the camera fallback chain}

\begin{codecard}
\lstinputlisting[firstline=257, lastline=289, firstnumber=257]{../hand_tracker.py}
\end{codecard}

\begin{itemize}
  \item \textbf{\code{save\_screenshot}} (lines 257--265) creates the directory if needed --- the
        \code{exist\_ok=True} flag makes repeated presses idempotent rather than errors --- and builds
        a timestamped filename. The timestamp uses \code{\%f} microseconds truncated to milliseconds
        (\code{[:-3]}), so two presses within the same second cannot overwrite each other.
  \item The \code{if cv2.imwrite(...)} test is not paranoia: \code{imwrite} returns \code{False} on
        failure (a full disk, a permission problem, an unsupported extension) rather than raising, so
        an unchecked call would fail silently and the user would simply never find the file.
  \item \textbf{\code{open\_camera}} (lines 268--289) is the fallback chain of
        Section~\ref{sec:camera}, covered in detail there. Two details are worth re-reading in the
        code: the backend list is chosen by \code{os.name == "nt"}, and the candidate list is built
        by filtering the requested index out of \code{(1, 0)} so no index is ever tried twice.
\end{itemize}

\subsection{The live capture loop}

\begin{codecard}
\lstinputlisting[firstline=292, lastline=338, firstnumber=292]{../hand_tracker.py}
\end{codecard}

\begin{itemize}
  \item \textbf{Lines 294--296}: the camera is opened \emph{first}, and a \code{None} return exits
        with status 1 before any model is constructed. Failing fast here means a missing webcam never
        pays the several-hundred-millisecond cost of loading the MediaPipe graph.
  \item \textbf{Lines 304--305}: a \code{deque(maxlen=30)} of instantaneous frame rates and a
        timestamp. The window gives a one-second moving average at 30 FPS --- long enough to smooth
        the reading, short enough to react to a real slowdown.
  \item \textbf{Lines 314--316, the per-frame core.} Flip, detect, annotate. Three lines, and
        \code{annotate\_hands} is shared with \code{--image} mode.
  \item \textbf{Lines 318--323, the FPS calculation.} This is the \code{time.time()} delta approach
        the specification requires: measure the interval since the previous frame, invert it for an
        instantaneous rate, and average the last thirty. The \code{if delta > 0} guard protects
        against a clock with insufficient resolution returning the same timestamp twice --- which
        would otherwise divide by zero.
  \item \textbf{Lines 325--326}: the HUD is drawn and the frame is shown. The order matters ---
        \code{imshow} displays whatever is in the buffer at that moment, so drawing after it would
        show a frame-delayed overlay.
  \item \textbf{Lines 328--332, input handling.} \code{cv2.waitKey(1)} waits one millisecond for a
        keypress, which also pumps the GUI event loop --- without it the window would freeze.
        \code{\& 0xFF} masks off platform-specific bits above the ASCII range, so the comparison
        against \code{ord("q")} works on every platform. Key 27 is \code{Esc}.
  \item \textbf{Lines 333--336, the \code{finally} block}: camera released, model closed, windows
        destroyed --- on every exit path, as required by Section~\ref{sec:teardown}.
\end{itemize}

\subsection{Headless image mode}

\begin{codecard}
\lstinputlisting[firstline=341, lastline=373, firstnumber=341]{../hand_tracker.py}
\end{codecard}

\begin{itemize}
  \item \textbf{Lines 342--346}: \code{cv2.imread} returns \code{None} for a missing or unreadable
        file --- it does not raise. Checking for \code{None} and returning 1 is what makes
        \code{--image does-not-exist.png} a clean failure that a CI script can detect from the exit
        code.
  \item \textbf{Lines 348--350}: resizing happens only when \emph{both} dimensions were supplied.
        The default of ``native resolution'' is the right choice for stills, where there is no
        frame-rate argument for downscaling.
  \item \textbf{Lines 352--357}: the tracker is constructed with \code{mode=True}
        (\code{static\_image\_mode}), because a single frame has no successor to track into.
  \item \textbf{Lines 359--363, the shared pipeline.} \code{find\_hands} $\to$
        \code{annotate\_hands} $\to$ \code{draw\_hud} --- the identical functions the live loop uses.
        The \code{try/finally} closes the model even if detection raises.
  \item \textbf{Lines 365--373, the output path.} \code{os.path.splitext} splits the input into stem
        and extension; the stem plus \code{\_out.png} is written \emph{next to the input}, so running
        the mode twice is idempotent and no directory configuration is needed. The reported
        \code{hand\_count()} comes from the cache built during \code{find\_hands}, which is why the
        log line can state a hand count without re-running detection.
\end{itemize}

\subsection{The command-line interface and entry point}

\begin{codecard}
\lstinputlisting[firstline=376, lastline=422, firstnumber=376]{../hand_tracker.py}
\end{codecard}

\begin{itemize}
  \item \textbf{Lines 377--381}: \code{argparse} is configured with a program name and a description,
        so \code{--help} produces usable output rather than a bare usage string.
  \item \textbf{Lines 382--384, the camera and size flags.} \code{--width} and \code{--height}
        default to \code{None}, not to 640 and 480. That deliberate choice is what allows
        \code{--image} mode to distinguish ``the user asked for a resize'' from ``the user accepted a
        default''; the live loop supplies 640$\times$480 itself when the flags are absent.
  \item \textbf{Line 388, \code{--thumb-mode}} uses \code{choices=("distance", "x")}. argparse then
        rejects a typo such as \code{--thumb-mode distancey} with a clear message, rather than
        silently falling back to the default.
  \item \textbf{Line 403, \code{--image}}: the flag that makes the whole application testable in CI,
        described in Section~\ref{sec:headless}.
  \item \textbf{Line 412--419, \code{main}}. Logging is configured from the parsed level, and the
        function returns an integer --- 0 or 1 --- rather than calling \code{sys.exit} itself. That
        makes \code{main()} callable from a test with an argument list, which is how a CLI is
        normally tested.
  \item \textbf{Lines 421--422, the entry point.} \code{raise SystemExit(main())} converts the return
        value into the process exit code, and the \code{\_\_name\_\_} guard ensures that importing
        the module (as the test suite does) never starts a camera.
\end{itemize}

\section{\code{tests/test\_hand\_logic.py} --- the offline test suite}
\label{sec:src-tests}

\subsection{The synthetic hands}

The file's first seventy lines define two coordinate sets by hand: \code{OPEN\_HAND} (all five
fingers extended) and \code{FIST} (all five folded). They are ordinary lists of $(x, y)$ tuples with
a comment naming each landmark. Because they are literals rather than captured data, every assertion
below is exact --- there is no tolerance and no randomness.

Four helpers manipulate them:

\begin{codecard}
\lstinputlisting[firstline=80, lastline=113, firstnumber=80]{../tests/test_hand_logic.py}
\end{codecard}

\begin{itemize}
  \item \code{\_with\_finger} and \code{\_with\_thumb} (lines 80--89) return a modified \emph{copy}
        of a hand with one finger, or the whole thumb, moved to a chosen position. Copying is
        essential: mutating the shared literal would make the tests order-dependent, so a test that
        passed alone could fail when run after another.
  \item \code{\_rotate} (lines 92--104) applies a real rotation matrix about the point set's
        centroid, rounding the result back to integers. It is the machinery behind the project's most
        important regression test --- the one that encodes the reviewer's original complaint about
        the horizontal thumb rule.
  \item \code{\_mirror\_x} (lines 107--113) negates every $x$ coordinate. Geometrically that is
        ``show the other side of the hand'', and it is exactly the transformation that must flip the
        palm normal's sign.
\end{itemize}

\subsection{The rotation-invariance test --- the project's key regression}

\begin{codecard}
\lstinputlisting[firstline=173, lastline=177, firstnumber=173]{../tests/test_hand_logic.py}
\end{codecard}

This four-line test is the reason the default thumb rule exists. It rotates the same open hand
through four angles and asserts that the thumb state never changes. The old horizontal rule fails
this test at 90\textdegree\ and 180\textdegree; the distance rule passes at every angle, for the
algebraic reason proved in Section~\ref{sec:thumb}. Turning a reviewer's complaint into an executable
assertion is what stops the bug from ever returning silently.

\subsection{The safe-result contract}

\begin{codecard}
\lstinputlisting[firstline=186, lastline=197, firstnumber=186]{../tests/test_hand_logic.py}
\end{codecard}

Four malformed inputs --- \code{None}, an empty list, a list one element too short, and a list of 21
\code{None} values --- all asserted to return exactly \code{[False] * 5}. The final line re-tests a
valid hand, guarding against the possibility that a guard was written so aggressively that it
suppresses correct input too. That is the shape of a complete boundary test: every invalid case
\emph{and} one valid case in the same test, so a regression in either direction fails.

\subsection{The fake MediaPipe results}

\begin{codecard}
\lstinputlisting[firstline=298, lastline=330, firstnumber=298]{../tests/test_hand_logic.py}
\end{codecard}

These five tiny classes are the reason \code{parse\_hand\_data} is testable without MediaPipe
installed. Each mimics exactly one level of the real result object's attribute chain:
\code{FakeResults} holds the two lists, \code{FakeHand} wraps a landmark list,
\code{FakeClassificationList} wraps a classification, and \code{FakeClassification} carries the
label. The full dotted chain --- from \code{multi\_handedness} through \code{classification[0]} down
to \code{label} --- therefore resolves exactly as it does in production, and the duck-typed parser
cannot tell the difference.

\code{FakeResults} has one particularly useful parameter:

\begin{itemize}
  \item \code{include\_handedness=False} omits the attribute \emph{entirely} on construction, which
        is how the ``missing \code{multi\_handedness}'' branch is covered. Passing an empty list
        would test something different: an attribute that exists but is empty. Both are real
        scenarios, and both are tested.
\end{itemize}

\subsection{Parse tests}

\begin{codecard}
\lstinputlisting[firstline=333, lastline=363, firstnumber=333]{../tests/test_hand_logic.py}
\end{codecard}

\begin{itemize}
  \item \textbf{\code{test\_parse\_two\_hands}} asserts on the label list, the landmark counts, the
        three-element \code{norm} tuples, the sign of \code{palm\_z}, and the orientation string ---
        then compares every pixel coordinate against the source with a $\pm 1$ tolerance. The
        tolerance is not sloppiness: the fake landmarks are built by dividing by 1000 and the parser
        multiplies by 1000 and truncates, so a half-unit difference is a floating-point artefact of
        the round trip, not a logic error. A genuinely wrong scale factor would be off by far more
        than one pixel.
  \item \textbf{\code{test\_parse\_zero\_hands}} covers the empty-but-present lists.
  \item \textbf{\code{test\_parse\_none\_and\_empty\_objects}} covers three distinct degenerate
        inputs: \code{None}, a bare \code{object()} with no attributes at all, and a result with
        neither list. All three must return \code{[]} rather than raising --- the ``never raises''
        contract is what allows the render loop to be written without a defensive \code{try} around
        it.
\end{itemize}

\begin{codecard}
\lstinputlisting[firstline=387, lastline=400, firstnumber=387]{../tests/test_hand_logic.py}
\end{codecard}

\begin{itemize}
  \item \textbf{\code{test\_parse\_skips\_incomplete\_hand}} builds a two-hand result where the first
        hand has only ten landmarks, and asserts that exactly one hand survives --- with the
        \emph{second} hand's label, proving the skip did not shift the handedness correspondence.
  \item \textbf{\code{test\_parse\_returns\_mirrored\_and\_back\_facing\_cases}} is the end-to-end
        confirmation of Section~\ref{sec:palm}: the same hand, mirrored in $x$, must flip the
        orientation label. It ties the geometry, the parser and the sign convention together in one
        assertion.
\end{itemize}

\begin{keyidea}
\textbf{Chapter summary.} The pure module converts coordinates once, applies three geometric rules
and two temporal filters, and never touches hardware or raises on bad input. The application layer
owns the model, the camera fallback chain, the HUD and the CLI, and routes both its live and headless
paths through the same annotation and drawing functions. The tests encode the two hardest-won
lessons of the project --- rotation invariance for the thumb, and a total, exception-free contract
for malformed landmarks --- as executable assertions rather than comments.
\end{keyidea}
